\chapter{Birth Control}

\section*{Definition and Overview}
Birth control, also called contraception, is any method used to prevent pregnancy. Options include hormonal methods such as the pill, implant, injection, and hormonal intrauterine devices; non-hormonal methods such as condoms, the copper intrauterine device, and diaphragms; fertility-awareness methods; and permanent methods such as sterilisation. Emergency contraception can prevent pregnancy after unprotected sex or a contraceptive failure.

The choice of method depends on effectiveness, health history, side effects, cost, and the length of time for which contraception is needed. Because no single method suits everyone, people are encouraged to discuss their options with a clinician and to review their choice as circumstances change.

\section*{Epidemiology}
Most sexually active people use some form of contraception during their reproductive years.

\begin{itemize}
    \item \textbf{Long-acting reversible methods:} the contraceptive implant and intrauterine devices are among the most effective methods and are increasingly used
    \item \textbf{Barrier methods:} condoms remain the most common method among young people and are the only method that also protects against sexually transmitted infections
    \item \textbf{Variation in use:} the choice of method varies with age, culture, access to services, and cost
    \item \textbf{Unplanned pregnancy:} remains common, and many occur when contraception is used inconsistently or not at all
\end{itemize}

\section*{Aetiology and Pathophysiology}
Contraceptive methods work by interfering with different steps in the process of conception.

\begin{itemize}
    \item \textbf{Combined hormonal methods:} the pill, patch, and vaginal ring use oestrogen and progestogen to stop ovulation and thicken cervical mucus
    \item \textbf{Progestogen-only methods:} the mini-pill, implant, injection, and hormonal intrauterine device thicken mucus, thin the uterine lining, and often suppress ovulation
    \item \textbf{Copper intrauterine device:} copper ions are toxic to sperm and interfere with fertilisation
    \item \textbf{Barrier methods:} condoms, diaphragms, and cervical caps physically block sperm from reaching the egg
    \item \textbf{Fertility-awareness methods:} identify the fertile window and avoid sex or use a barrier method during it
    \item \textbf{Sterilisation:} blocks or cuts the tubes so that sperm and egg cannot meet, either a vasectomy in men or tubal occlusion in women
\end{itemize}

\section*{Clinical Features}
There is no single set of symptoms; the features relate to the method chosen, how it is used, and its effects on the body.

\begin{itemize}
    \item \textbf{Combined hormonal methods:} many users have regular, lighter, and less painful periods
    \item \textbf{Progestogen-only methods:} can cause irregular or absent bleeding, with lighter periods over time
    \item \textbf{Hormonal intrauterine devices:} reduce menstrual bleeding, sometimes stopping periods altogether
    \item \textbf{Copper intrauterine device:} may cause heavier and more painful periods
    \item \textbf{Barrier methods:} have no hormonal effects and are used only at the time of sex
    \item \textbf{Fertility-awareness methods:} require daily tracking of menstrual cycle signs
    \item \textbf{Possible side effects:} breast tenderness, mood changes, nausea, and weight change, depending on the method
\end{itemize}

\section*{Differential Diagnosis}
Before starting any method, a clinician will rule out conditions that affect the choice of contraception.

\begin{itemize}
    \item Current or possible pregnancy before starting a method
    \item Undiagnosed abnormal vaginal bleeding
    \item High blood pressure, migraine with aura, or a personal or family history of blood clots, which are relevant to combined hormonal methods
    \item Liver disease or certain cancers, which may limit hormonal contraception
    \item Pelvic infection, which must be treated before intrauterine device insertion
    \item Medicines or conditions that can reduce the effectiveness of hormonal methods
\end{itemize}

\section*{When to Seek Medical Care}
See a doctor, nurse, or family planning clinic for help choosing or starting contraception, for a method that is not working, or if you think you may be pregnant. Situations such as a condom slipping or breaking, a missed contraceptive pill, or a late injection require prompt advice about emergency contraception.

\textbf{Call 000 (or local emergency services) for:}
\begin{itemize}
    \item Severe lower abdominal pain with a positive pregnancy test, which may indicate an ectopic pregnancy
    \item Sudden breathlessness or chest pain after starting a new hormonal method
    \item Heavy bleeding, severe pain, or fainting after a procedure such as intrauterine device insertion or sterilisation
\end{itemize}

\section*{Diagnosis and Investigations}
Before prescribing, clinicians assess the suitability of each method.

\begin{itemize}
    \item \textbf{History:} age, smoking status, blood pressure, migraine history, menstrual history, and personal or family history of blood clots or breast cancer
    \item \textbf{Pregnancy test:} to confirm the person is not pregnant before starting a method
    \item \textbf{Blood pressure measurement:} required before combined hormonal methods
    \item \textbf{Body mass index calculation:} some methods are less effective or carry greater risk at high body weight
    \item \textbf{STI screening:} recommended before intrauterine device insertion
    \item \textbf{Pelvic examination:} may be performed before some methods
\end{itemize}

\section*{Management}
Methods are chosen with a clinician and reviewed over time.

\begin{itemize}
    \item \textbf{Combined oral contraceptive pill, patch, or ring:} highly effective when taken as directed
    \item \textbf{Progestogen-only pill (mini-pill):} an option for those who cannot take oestrogen
    \item \textbf{Contraceptive implant:} a small rod placed in the upper arm, effective for several years
    \item \textbf{Contraceptive injection:} a progestogen injection given every few months
    \item \textbf{Intrauterine devices:} copper or hormonal devices inserted by a trained clinician, effective for 5 to 10 years
    \item \textbf{Barrier methods:} male and female condoms, diaphragms, and cervical caps
    \item \textbf{Fertility awareness:} requires training and a strict daily routine
    \item \textbf{Permanent methods:} vasectomy and tubal occlusion for those who are certain they want no more children
    \item \textbf{Emergency contraception:} a morning-after pill or a copper intrauterine device, available up to five days after unprotected sex
\end{itemize}

\section*{Complications}
\begin{itemize}
    \item Combined hormonal methods slightly raise the risk of blood clots, stroke, and heart attack, particularly in smokers and older women
    \item Some hormonal methods can raise blood pressure and worsen migraine
    \item Intrauterine device insertion can cause pain and cramping and, rarely, perforation of the uterus
    \item The copper intrauterine device may cause heavier periods and, in some women, anaemia
    \item Implants and injections often cause irregular bleeding, which leads many users to discontinue them
    \item Barrier methods can fail with incorrect use, and condoms can slip or tear
    \item Sterilisation is intended to be permanent, and reversal is difficult
    \item Contraception does not protect against sexually transmitted infections unless condoms are used
\end{itemize}

\section*{Prognosis and Outcomes}
With correct and consistent use, most methods are highly effective at preventing pregnancy. Long-acting reversible methods have the lowest failure rates because they do not depend on daily use.

When used perfectly, the pill, implant, and intrauterine devices are all more than 99 per cent effective; typical-use failure rates are higher, mainly because people forget to use them correctly. Fertility returns promptly after stopping most methods, although the injection can delay the return of ovulation for many months. There is no evidence that any contraceptive method causes lasting infertility.

\section*{Prevention and Self-Care}
\begin{itemize}
    \item Choose a method that fits your lifestyle; long-acting methods suit people who may forget a daily pill
    \item Use condoms as well as hormonal contraception to reduce the risk of sexually transmitted infections
    \item Take the pill at the same time each day and follow the instructions for missed pills
    \item Keep a reminder of injection and implant due dates
    \item Do not use combined hormonal contraception if you smoke and are over 35, unless advised otherwise by a doctor
    \item Seek emergency contraception within the recommended window after unprotected sex
    \item Discuss side effects with your clinician rather than stopping the method without advice
    \item Attend regular check-ups, including blood pressure checks and cervical screening as recommended
\end{itemize}

\section*{Sources and Further Reading}
The following organisations provide reliable information on birth control for patients, families, and health professionals.

\begin{itemize}
    \item NHS. \textit{Information on contraception and available birth control methods.} https://www.nhs.uk/
    \item Mayo Clinic. \textit{Overview of birth control options and their effectiveness.} https://www.mayoclinic.org/
    \item MedlinePlus. \textit{Health information on birth control and contraception for patients.} https://medlineplus.gov/
    \item World Health Organization. \textit{Resources on family planning and contraception.} https://www.who.int/
    \item Centers for Disease Control and Prevention. \textit{Resources on contraception and family planning.} https://www.cdc.gov/
\end{itemize}
